% 03_introduction.tex — Introducción (Sección 1), estructura IMRAD:
%   1.1 Contexto → 1.2 Problema de control → 1.3 Estado del arte en dominios
%   vecinos → 1.4 Brecha y justificación → 1.5 Objetivo y preguntas.
% La defensa del diseño (revisión de alcance vs. narrativa/agnóstica) se movió
% a §2.1 (Métodos), donde corresponde en IMRAD.
\setcounter{section}{0}
\renewcommand{\thesection}{\arabic{section}}
\setlength{\parindent}{15pt}

\section{Introducción}

\subsection{Contexto industrial: el pH en la clarificación azucarera}

La caña de azúcar aporta la mayor parte de la sacarosa comercializada en
el mundo \parencite{rein2007cane}. Dentro de la cadena de fabricación
---preparación de caña, extracción, \emph{clarificación y encalado},
evaporación, cristalización, centrifugación y secado---, la clarificación
es la primera operación de purificación del jugo crudo y la más sensible
al control de pH. La adición controlada de lechada de cal
(\textit{liming}) eleva el pH desde su valor natural ácido
($\approx 4.8$--$5.8$) hasta el rango objetivo de alcalinización
($\approx 7.0$--$7.2$), lo que precipita coloides, ceras, materia
colorante y sales inorgánicas y protege la sacarosa de la hidrólisis ácida
\parencite{rein2007cane}.

Las desviaciones del \textit{setpoint} tienen consecuencias económicas en
ambos sentidos. Si el pH desciende por debajo de $\approx 6.8$, la acidez
residual cataliza la inversión de sacarosa en glucosa y fructosa y reduce
el rendimiento azucarero; si excede $\approx 7.4$, el exceso de cal
favorece la formación de sacaratos de calcio y compuestos coloreados que
afectan la calidad del azúcar y la eficiencia de los evaporadores
\parencite{shinskey1996process,rein2007cane}. Por ello el encalado se
reconoce como una operación cuya regulación tiene retorno directo sobre
la rentabilidad del ingenio.

\subsection{El problema de control del pH: ganancia variable y no linealidad}

El encalado es una instancia del problema clásico de neutralización
ácido--base: la curva de titulación es fuertemente sigmoidal, con una
ganancia estática que varía en varios órdenes de magnitud alrededor del
punto de equivalencia \parencite{shinskey1996process}, comportamiento que
se formaliza con los invariantes de reacción de
\textcite{gustafsson1983dynamic}. Un controlador PID sintonizado en un
punto de operación pierde desempeño al alejarse de él, con oscilaciones
y sobredosificación de reactivo. A esa no linealidad se superponen tres
dificultades atribuidas en la literatura tecnológica al jugo de caña
\parencite{rein2007cane,garcia1989model}: (i)~\textbf{variabilidad
temporal} del jugo crudo, cuya composición ácida depende de la variedad,
la madurez y la campaña; (ii)~\textbf{asimetría de instrumentación}, con
mediciones DCS a resolución de minutos y análisis de laboratorio cada
1--2~horas; y (iii)~\textbf{ensuciamiento} (\textit{fouling}) de los
electrodos por depósitos de cal, que introduce derivas y exige
recalibración frecuente. La Figura~\ref{fig:causal} resume esta cadena
causal, de la no linealidad fisicoquímica a las limitaciones del PID, y los
mecanismos mediante los cuales el aprendizaje automático (AA) y el control
avanzado, discutidos en \ref{sub:otros_dominios} y
\ref{sub:transferibilidad}, podrían mitigarlas.

Conviene precisar el alcance de esta premisa: que el PID se degrada en el
encalado es una consecuencia prevista por la teoría de control de pH y
reportada de forma cualitativa por los propios estudios del dominio
\parencite{murugesan2020computeraided}; no se localizó, en las fuentes
revisadas, un estudio que cuantifique el desempeño de lazos PID en
clarificadores industriales (error de pH, consumo de cal, pérdidas de
sacarosa). Esa ausencia de línea base empírica es, en sí misma, parte del
problema que esta revisión documenta.

\nolinenumbers
\begin{figure}[H]
\centering
\captionsetup{width=0.92\textwidth}
\caption{Esquema causal propio: de la no linealidad fisicoquímica del
encalado a las limitaciones del control PID, y mecanismos por los cuales
el aprendizaje automático y el control avanzado podrían mitigarlas.}
\label{fig:causal}
\resizebox{\textwidth}{!}{% diagrama_causal.tex — Esquema causal propio (Observación 5, revisión externa)
% No linealidad fisicoquímica del encalado -> limitaciones del PID ->
% mecanismos de mejora mediante AA/control avanzado.
% Se incluye vía \input{} dentro de un entorno figure en 03_introduction.tex.
% Layout por coordenadas absolutas (x,y) en cm, no posicionamiento relativo
% encadenado, para evitar colisiones entre columnas de distinta altura.
\begin{tikzpicture}[
  font=\sffamily\scriptsize,
  box/.style = {draw, rounded corners=2pt, align=center, text width=3.0cm,
                minimum height=1.6cm, inner sep=3pt, line width=0.6pt},
  causa/.style   = {box, fill=blue!6},
  efecto/.style  = {box, fill=red!6},
  mejora/.style  = {box, fill=green!6},
  flecha/.style  = {-{Stealth[length=2mm]}, line width=0.6pt},
  etiqueta/.style = {font=\sffamily\bfseries\scriptsize, text=black!65},
]

% ── Encabezados de columna ─────────────────────────────────────────────────────
\node[etiqueta, align=center, text width=2.6cm] at (1.6, 9.5)
  {No linealidad\\fisicoquímica\\del encalado};
\node[etiqueta, align=center, text width=3.0cm] at (6.7, 9.5)
  {Limitación del\\control PID};
\node[etiqueta, align=center, text width=2.8cm] at (11.8, 9.5)
  {Mecanismos de\\AA / control\\avanzado};

% ── Columna 1: causa raíz fisicoquímica (y = 7.5, 4.6, 1.5) ───────────────────
\node[causa] (sigmoidal) at (1.6, 7.5)   {Curva de titulación\\\textit{sigmoidal}:
  ganancia variable en\\varios órdenes de\\magnitud cerca de\\pH~$\approx$~7.0--7.2};
\node[causa] (variabilidad) at (1.6, 4.6) {Variabilidad del\\jugo crudo\\
  (variedad, madurez,\\campaña)};
\node[causa] (fouling) at (1.6, 1.5)   {\textit{Fouling} de\\electrodos y
  \\asimetría DCS~vs.\\laboratorio};

% ── Columna 2: nodo central PID (y=4.6) y sus dos consecuencias (y=1.1) ───────
\node[efecto] (pid)   at (6.7, 4.6) {PID sintonizado en\\un punto de
  operación\\pierde desempeño\\al alejarse de él};
\node[efecto, text width=3.2cm] (oscil) at (4.8, 1.0) {Oscilaciones\\sostenidas y
  \\sobredosificación\\de cal};
\node[efecto, text width=3.2cm] (costo) at (8.6, 1.0) {Costo económico\\doble:
  pH~$<$~6.8~$\to$\\inversión de sacarosa;\\pH~$>$~7.4~$\to$\\sacaratos de calcio};

% ── Columna 3: mecanismos de mejora (y = 7.5, 4.6, 1.5) ───────────────────────
\node[mejora] (fuzzy) at (11.8, 7.5)   {Control neuro-difuso\\orientado a la zona
  \\de máxima sensibilidad\\(punto de equivalencia)};
\node[mejora] (nmpc)  at (11.8, 4.6) {NMPC / caja gris con\\identificación de
  \\sistema explícita de\\la no linealidad};
\node[mejora] (soft)  at (11.8, 1.5) {Sensores virtuales /\\\textit{data-driven}
  \\robustos a \textit{fouling}\\y a baja frecuencia\\de muestreo};

% ── Flechas causales (columna 1 -> columna 2) ─────────────────────────────────
\draw[flecha] (sigmoidal) -- (pid);
\draw[flecha] (variabilidad) -- (pid);
\draw[flecha] (fouling) -- (pid);
\draw[flecha] (pid) -- (oscil);
\draw[flecha] (pid) -- (costo);

% ── Flechas causales (columna 1 -> columna 3) ─────────────────────────────────
\draw[flecha] (sigmoidal) -- (fuzzy);
\draw[flecha] (variabilidad) -- (nmpc);
\draw[flecha] (fouling) -- (soft);

% ── Flechas de mitigación (columna 3 -> columna 2, discontinuas) ──────────────
\draw[flecha, dashed, blue!60!black] (fuzzy) -- (pid.north east)
  node[midway, above, sloped, etiqueta, text=blue!60!black] {\scriptsize mitiga};
\draw[flecha, dashed, blue!60!black] (nmpc) -- (pid.east)
  node[midway, above, sloped, etiqueta, text=blue!60!black] {\scriptsize mitiga};
\draw[flecha, dashed, blue!60!black] (soft) -- (pid.south east)
  node[midway, below, sloped, etiqueta, text=blue!60!black] {\scriptsize mitiga};

\end{tikzpicture}
}
\begin{Nota}[Nota]{0.92\textwidth}{Elaboración propia. Las flechas sólidas
representan relaciones causales introducidas en los párrafos anteriores y
desarrolladas con su respaldo bibliográfico en \ref{sub:otros_dominios};
las flechas discontinuas señalan el mecanismo de mitigación propuesto por
cada familia de métodos, desarrollado en \ref{sub:transferibilidad} y en
la agenda de investigación (\ref{sub:agenda}). El esquema es conceptual y
no representa un diagrama de bloques de control formal.}
\end{Nota}
\end{figure}
\linenumbers

\subsection{Avances en AA y control avanzado en otros dominios}
\label{sub:otros_dominios}

Fuera de la industria azucarera, la dificultad del control de pH ha
motivado cuatro décadas de desarrollo metodológico. Los trabajos
fundacionales formularon el control adaptativo no lineal basado en
invariantes de reacción \parencite{henson1994adaptive}, el control
predictivo con identificación no lineal en tiempo real
\parencite{proll1994modelpredictive} y las redes neuronales para capturar
la relación entrada--pH sin un modelo mecanístico explícito
\parencite{palancar1998phcontrol}; en paralelo se propusieron esquemas
proporcional-integrales avanzados con inyección dual para mitigar la
ganancia variable \parencite{hurowitz2000control}.

Sobre esa base se ha consolidado un repertorio diverso. El control
predictivo lineal y no lineal se ha aplicado a la regulación de pH en
bioprocesos y fotobiorreactores
\parencite{berenguel2004model,otalora2026enhancing} y, más
recientemente, se ha combinado con identificación por aprendizaje
profundo para sintetizar controladores del proceso de neutralización
\parencite{durairaj2023deep}. Los enfoques difusos y neuro-difusos han
resultado adecuados para operar cerca del punto de equivalencia, donde la
sensibilidad del proceso es máxima
\parencite{herediamolinero2014feedback,suguna2026advanced}. En la frontera
actual, el aprendizaje por refuerzo ha permitido diseñar políticas
autónomas para la neutralización de efluentes de galvanoplastia y para
procesos multivariables de pH
\parencite{alvesgoulart2020autonomous,panjapornpon2022reinforcement}, y
los enfoques \textit{data-driven} y de sensores virtuales han abordado la
predicción de pH y la dosificación de reactivos cuando la
instrumentación física es limitada o poco fiable
\parencite{son2021deep,xu2022datadriven,bardeeniz2025artificial}.

Este cuerpo de conocimiento abarca el tratamiento de aguas residuales, los
fotobiorreactores, la desalación y la galvanoplastia, y su madurez se
apoya sobre todo en formulación y validación en simulación, con
implementaciones experimentales en algunos casos. Cabe preguntarse en
qué medida estos avances han llegado al proceso de clarificación del
jugo de caña, cuya dinámica de pH es análoga pero presenta
particularidades ---composición coloidal, multiplicidad de ácidos débiles,
estacionalidad y \textit{fouling}--- que podrían condicionar la
transferencia.

\subsection{Brecha de investigación y justificación del enfoque}
\label{sub:brecha}

Una búsqueda exploratoria sobre Scopus, Web of Science e IEEE~Xplore
sugería que la evidencia \textit{intra-dominio} (\textit{in-domain}) sobre
AA y control avanzado aplicados específicamente al pH del proceso
azucarero era escasa y no existía un mapeo sistemático de su extensión y
naturaleza. Ese contexto ---tema acotado, evidencia dispersa en tipos de
documento heterogéneos y objetivo de cartografiar más que de sintetizar
efectos--- corresponde a una \textbf{revisión de alcance}
\parencite{arksey2005scoping,levac2010scoping,peters2020scoping}. Se
adoptó un protocolo sistemático, y no una revisión narrativa, porque solo
una estrategia de búsqueda reproducible y un embudo de decisiones
auditable permiten distinguir entre ``no hay evidencia'' y ``no se buscó
exhaustivamente'', distinción central cuando el corpus elegible es
pequeño \parencite{lang2007empty,yaffe2012empty}. La justificación
completa del diseño se presenta en \ref{sub:diseno}; la síntesis
narrativa de la evidencia transferible, en \ref{sub:transferibilidad}.

\subsection{Objetivo y preguntas de investigación}
\label{sub:rqs}

El \textbf{objetivo} es cartografiar la extensión, la naturaleza y el
nivel de madurez de la evidencia sobre AA, control avanzado y sensores
virtuales aplicados a la regulación del pH en la clarificación y el
encalado del jugo azucarero, e identificar vacíos que orienten una agenda
de investigación. Bajo un marco PICOC adaptado
(\textit{Population}: procesos de clarificación y encalado de jugo de caña
o remolacha; \textit{Intervention}: técnicas de AA, control avanzado o
sensores virtuales; \textit{Comparison}: PID, dosificación manual o lazo
abierto; \textit{Outcomes}: desempeño del control de pH y evidencia de
modelado; \textit{Context}: revistas indexadas en SJR, cualquier cuartil,
2015--2026) se formulan tres preguntas con correspondencia biunívoca con
las consultas de búsqueda:

\begin{description}[style=nextline,leftmargin=2cm,labelindent=0cm]
\item[\textbf{RQ1}] ¿Cómo se ha modelado y estimado la dinámica del pH
en la clarificación azucarera mediante sensores virtuales, identificación
de sistemas y modelos \textit{data-driven}?
\item[\textbf{RQ2}] ¿Qué estrategias de control avanzado (MPC/NMPC,
lógica difusa, metaheurísticas) se han aplicado al lazo de pH y cómo se
comparan con el control PID convencional?
\item[\textbf{RQ3}] ¿Qué técnicas de aprendizaje automático, aprendizaje
por refuerzo y modelos híbridos se han empleado, y qué barreras limitan
su adopción industrial?
\end{description}

\noindent\textbf{Contribución.} Hasta donde alcanza el conocimiento de los
autores, este es el primer mapeo sistemático, reportado bajo PRISMA-ScR,
de la intersección \{AA/control avanzado\,$\times$\,pH\,$\times$\,industria
azucarera\}. Aporta: \emph{(i)}~un corpus auditado registro a registro,
con fiabilidad de cribado cuantificada y análisis de sensibilidad de los
criterios de elegibilidad; \emph{(ii)}~la caracterización del campo
adyacente mediante bibliometría y una verificación léxica independiente;
y \emph{(iii)}~una agenda priorizada por transferibilidad técnica y nivel
de madurez, orientada a cerrar la distancia entre simulación y validación
en planta.
